\appendix
\chapter{Appendix}
\begin{spacing}{1.0}
\small

\section{Exploratory Plots}
\label{app:exploratory-plots}

This appendix collects the supplementary exploratory figures referenced in the
data chapter.

\begin{figure}[p]
    \centering
    \includegraphics[
        width=0.9\textwidth,
        height=0.82\textheight,
        keepaspectratio
    ]{skewness_plots.png}
    \caption{Raw versus $\log(1+x)$ distributions for the heavily skewed
    numerical features identified during exploratory analysis.}
    \label{fig:skewness-log1p}
\end{figure}

\begin{figure}[p]
    \centering
    \includegraphics[
        width=\textwidth,
        height=0.78\textheight,
        keepaspectratio
    ]{categorical_plots.png}
    \caption{Frequency distributions of the bucketed categorical features used
    in the final preprocessing pipeline.}
    \label{fig:categorical-buckets}
\end{figure}

\begin{figure}[p]
    \centering
    \includegraphics[
        width=\textwidth,
        height=0.78\textheight,
        keepaspectratio
    ]{correlation-heatmap.png}
    \caption{Correlation heatmap over the numerical features used during
    exploratory analysis.}
    \label{fig:correlation-heatmap}
\end{figure}

\section{Per-Attack Category Tables}
\label{app:per-attack-plots}

The following figures reproduce the per-attack-category table outputs
reported in each model notebook.

Per-attack precision is uniformly 1.0 by construction of this breakdown.
Each table groups only true attack rows after excluding Normal traffic, so
no false-positive cases remain within an attack-category subset. As a
result, recall and F1 are the more informative per-category metrics.

\begin{figure}[p]
    \centering
    \includegraphics[
        width=\textwidth,
        height=0.72\textheight,
        keepaspectratio
    ]{logistic_per_attack_table.png}
    \caption{Logistic Regression per-attack-category table output.}
    \label{fig:logistic-per-attack}
\end{figure}

\begin{figure}[p]
    \centering
    \includegraphics[
        width=\textwidth,
        height=0.72\textheight,
        keepaspectratio
    ]{mlp_per_attack_table.png}
    \caption{MLP per-attack-category table output.}
    \label{fig:mlp-per-attack}
\end{figure}

\begin{figure}[p]
    \centering
    \includegraphics[
        width=\textwidth,
        height=0.72\textheight,
        keepaspectratio
    ]{random_forest_per_attack_table.png}
    \caption{Random Forest per-attack-category table output.}
    \label{fig:rf-per-attack}
\end{figure}

\begin{figure}[p]
    \centering
    \includegraphics[
        width=\textwidth,
        height=0.72\textheight,
        keepaspectratio
    ]{xgboost_per_attack_table.png}
    \caption{XGBoost per-attack-category table output.}
    \label{fig:xgb-per-attack}
\end{figure}

\section{Feature Breakdown}
\label{app:feature-breakdown}

\begin{table}[htbp]
    \centering
    \caption{Raw categorical and numerical features in the original dataset.}
    \label{tab:feature-breakdown}
    \begin{tabularx}{\linewidth}{>{\raggedright\arraybackslash}p{0.2\linewidth} X}
        \toprule
        Feature Type & Variables \\
        \midrule
        Categorical (3) &
        \texttt{proto}, \texttt{service}, \texttt{state} \\
        Numerical (37) &
        \texttt{ackdat}, \texttt{ct\_dst\_ltm}, \texttt{ct\_dst\_sport\_ltm},
        \texttt{ct\_dst\_src\_ltm}, \texttt{ct\_flw\_http\_mthd},
        \texttt{ct\_ftp\_cmd}, \texttt{ct\_src\_dport\_ltm},
        \texttt{ct\_src\_ltm}, \texttt{ct\_srv\_dst}, \texttt{ct\_srv\_src},
        \texttt{ct\_state\_ttl}, \texttt{dbytes}, \texttt{dinpkt},
        \texttt{djit}, \texttt{dload}, \texttt{dloss}, \texttt{dmean},
        \texttt{dpkts}, \texttt{dttl}, \texttt{dur}, \texttt{dwin},
        \texttt{is\_ftp\_login}, \texttt{is\_sm\_ips\_ports}, \texttt{rate},
        \texttt{response\_body\_len}, \texttt{sbytes}, \texttt{sinpkt},
        \texttt{sjit}, \texttt{sload}, \texttt{sloss}, \texttt{smean},
        \texttt{spkts}, \texttt{sttl}, \texttt{swin}, \texttt{synack},
        \texttt{tcprtt}, \texttt{trans\_depth} \\
        \bottomrule
    \end{tabularx}
\end{table}

\section{Categorical Frequency Tables}
\label{app:categorical-frequency}

\begin{table}[htbp]
    \centering
    \caption{Retained categorical levels based on training-split frequencies.}
    \label{tab:categorical-retained}
    \begin{tabularx}{\linewidth}{>{\raggedright\arraybackslash}p{0.14\linewidth} X >{\raggedright\arraybackslash}p{0.12\linewidth}}
        \toprule
        Feature & Retained levels (counts) & Collapsed into \texttt{other} \\
        \midrule
        \texttt{proto} &
        \texttt{tcp} (43{,}095), \texttt{udp} (29{,}418),
        \texttt{unas} (3{,}515), \texttt{arp} (987), \texttt{ospf} (676),
        \texttt{sctp} (324) &
        4{,}317 \\
        \texttt{service} &
        \texttt{-} (47{,}153), \texttt{dns} (21{,}367),
        \texttt{http} (8{,}287), \texttt{smtp} (1{,}851),
        \texttt{ftp} (1{,}552), \texttt{ftp-data} (1{,}396),
        \texttt{pop3} (423), \texttt{ssh} (204) &
        99 \\
        \texttt{state} &
        \texttt{FIN} (39{,}339), \texttt{INT} (34{,}163),
        \texttt{CON} (6{,}982), \texttt{REQ} (1{,}842) &
        6 \\
        \bottomrule
    \end{tabularx}
\end{table}

\end{spacing}
\newpage
